\begin{mainquestion}{diodeconnected}{The Terrific Two Terminal Transistor}

\questionpoints
The circuit shown in Figure (a) below converts the MOSFET into a two-terminal device by shorting the gate and drain.

\begin{center}
\begin{circuitikz}[american]
    % (a) Two-terminal configuration
    
    % Place the NMOS transistor
    \draw (0,0) node[nmos] (Q1) {};
    
    % Drain (Top) Wire
    \draw (Q1.D) ++(0,1.5) 
        node[above, yshift=3pt] {$+$}      % Polarity +
        to[short, o-, i^=$i$] ++(0,-1)      % Current 'i' on the right side
        coordinate(junction) 
        -- (Q1.D);

    % Gate Connection
    \draw (Q1.G) |- (junction) node[circ] {};
    
    % Source (Bottom) Wire
    \draw (Q1.S) -- ++(0,-1) 
        node[below, yshift=-3pt] {$-$}      % Polarity -
        to[short, -o] ++(0,0);
    
    % Voltage Label 'v'
    \draw (Q1.D) ++(2.0, 1.5) to[open, v^=$v$] ++(0, -3.5);
    
    % Figure Label
    \draw (0,-2.5) node {(a) Two-terminal configuration};

\end{circuitikz}
\end{center}

% =========================================================
\begin{partquestion}{derive}
{Derive an expression for the current versus voltage ($i$ vs. $v$) of this NMOS transistor circuit for both cutoff and non-cutoff operation. What operating regime is it in? You may assume this is an enhancement-mode transistor, i.e., $V_t > 0$.}

\begin{answer}
        Here, $v_{GS} = v_{DS} = v$. Since $v_{DS} > v_{GS} - V_t$ (because $v > v - V_t$ implies $0 > -V_t$), the device will always be in \textbf{saturation} provided $v > V_t$.
        
        The equations are:
        \[
        i = 
        \begin{cases} 
        0 & \text{for } v < V_t \quad (\text{Cutoff}) \\
        \frac{1}{2} k_n' \frac{W}{L} (v - V_t)^2 & \text{for } v \geq V_t \quad (\text{Saturation})
        \end{cases}
        \]
\end{answer}

\end{partquestion}

% =========================================================
\begin{partquestion}{explain}
{This configuration is called a \textit{diode-connected transistor}. Why do you think it has such a name? Sketch or plot the $i$-$v$ characteristic curve to support your answer. How does the current growth compare to that of the Shockley model?}

\begin{answer}
        Like a diode, this diode-connected transistor acts as a switch at a certain applied voltage (in this case $V_t$) below which the device is off. Above this voltage, the current increases as the applied voltage increases. 
        
        \begin{center}
        \begin{tikzpicture}
            \draw[->] (-1,0) -- (4,0) node[right] {$v$};
            \draw[->] (0,-0.5) -- (0,4) node[above] {$i$};
            \draw[thick, blue] (0,0) -- (1.5,0) node[below, black] {$V_t$} .. controls (2.5,0.5) and (3,2) .. (3.5,3.5);
        \end{tikzpicture}
        \end{center}
        
        However, the growth here is \textbf{quadratic}, which is not as fast as the exponential growth seen in the Shockley diode model.
\end{answer}

\pagebreak
\end{partquestion}

% =========================================================
\begin{partquestion}{biascircuit}
{Now consider this diode-connected transistor is used in the bias circuit shown in Figure (b). The supply voltage is $V_{DD} = +3 \text{ V}$ and the resistor is $R = 2 \text{ k}\Omega$.
    
    Assume $V_t = 1.0 \text{ V}$ and $k_n' (W/L) = 1.0 \text{ mA/V}^2$. Calculate $V_{out}$.
    
    \begin{center}
    \begin{circuitikz}[american]
        % (b) Resistive load circuit
        
        % VDD and Resistor
        \draw (0,4) node[vcc] {$V_{DD} = +3\text{V}$}
        to[R, l=$R$, i_=$I_D$] (0,2) 
        coordinate(out);
        
        % Output terminal
        \draw (out) to[short, -o] ++(1.5,0) node[right] {$V_{out}$};
        
        % Transistor Q1
        \draw (0,0) node[nmos] (Q1) {};
        \draw (out) -- (Q1.D); 
        \draw (Q1.S) node[ground]{}; 
        
        % Diode connection
        \draw (Q1.G) -- ++(0,0.5) -| (Q1.D) node[circ, pos=1]{};
        
        % Figure Label
        \draw (0,-2) node {(b) Resistive load circuit};
    \end{circuitikz}
    \end{center}}

\begin{answer}
        From KVL: $V_{DD} - I_D R = V_{out}$.
        Substituting saturation current:
        \[ \frac{3 - V_{out}}{2} = \frac{1}{2} (1) (V_{out} - 1)^2 \]
        \[ 3 - V_{out} = (V_{out} - 1)^2 \implies V_{out}^2 - V_{out} - 2 = 0 \]
        \[ (V_{out} - 2)(V_{out} + 1) = 0 \]
        Since $V_{out} > V_t$, the solution is $\mathbf{V_{out} = 2.0 \text{ V}}$.
\end{answer}

\end{partquestion}

\end{mainquestion}
